% !TEX root = ../main.tex

% Draft section for later inclusion in the paper.
% Required preamble packages:
%   \usepackage{algorithm}
%   \usepackage{algpseudocode}

\section{Closing the Market}
\label{sec:market-close}

In this appendix, we describe Algorithm~\ref{alg:market-close} for closing a batch auction supporting three order types: limit-on-close (LOC), market-on-close (MOC), and imbalance-only (IO) orders. Real-world exchanges differ in the precise specification of their closing procedures. We therefore present a stylized, ``textbook'' algorithm: it begins with the call-auction procedure described by Harris~\cite{Har03} and incorporates selected features of the Nasdaq Closing Cross.

The algorithm selects the clearing price in three phases using the ordinary LOC and MOC orders. First, it identifies the price or prices that maximize executable volume. If a unique price maximizes volume, that price is selected. Otherwise, the volume-maximizing prices form a \emph{clearing-price interval}. Second, the algorithm restricts this interval to the prices that minimize the ordinary-order imbalance. If this produces a unique price, that price is selected. Third, if multiple prices remain, the algorithm selects a price from remaining interval according to some rule (\eg  midpoint, using a fixed rounding rule if the midpoint does not lie on an admissible price tick). After the clearing price has been selected, eligible IO orders are admitted only on the short side of the market and only to the extent that they offset the remaining ordinary-order imbalance. The resulting executions are then allocated at the selected uniform clearing price.

Our zk-FBA proves the output of Algorithm~\ref{alg:market-close} is correct. This makes one modification from the high-level description above. Instead of outputting a single market clearing price, it outputs the price interval from phase 1 and, if applicable, phase 2. Given the price interval from phase 2, anyone can recompute the single market clearing price implied by the interval. 

\subsection{Notation and Order Eligibility}
\label{sec:close-notation}

Let $\mathcal{O}^{\mathrm{ord}}$ contain the ordinary closing orders---the
limit-on-close (LOC) and market-on-close (MOC) orders---and let
$\mathcal{O}^{\mathrm{IO}}$ contain the imbalance-only (IO) orders. For an
order $o$, let $q(o)$ denote its quantity and $e(o)$ its entered price. The
admissible price grid is

\[
    \mathcal{P}=\{\pi_0<\pi_1<\cdots<\pi_{N-1}\}.
\]

For aggregation purposes, a buy MOC is entered at $\pi_{N-1}$ and a sell MOC
at $\pi_0$, making each eligible at every admissible price. A buy LOC with
limit $\ell$ is eligible at $\pi_i$ when $\pi_i\leq\ell$; a sell LOC is
eligible when $\pi_i\geq\ell$. IO orders obey the same price conditions, but
do not participate in selecting the clearing price and may execute only on the
short side after that price has been selected. A higher entered price is more
aggressive for a buy order, while a lower entered price is more aggressive for
a sell order.

The market quantities used by the closing procedure are:

\begin{description}
    \item[$B_i,A_i$] ordinary buy and sell quantities entered exactly at
        price $\pi_i$;
    \item[$D_i,S_i$] ordinary demand and supply eligible at $\pi_i$;
    \item[$V_i$] executable ordinary volume $\min\{D_i,S_i\}$;
    \item[$G_i$] signed imbalance $S_i-D_i$, so $G_i<0$ means demand is long
        and $G_i>0$ means supply is long;
    \item[$\Delta_i$] absolute imbalance $|G_i|$;
    \item[$V_{\max},\Delta_{\min}$] maximum executable ordinary volume and
        minimum imbalance among the volume-maximizing prices;
    \item[$i^\star,\pi^\star$] clearing-price index and clearing price, with
        $\pi^\star=\pi_{i^\star}$;
    \item[$P$] ordinary long-side quantity entered exactly at $\pi^\star$;
    \item[$J,I,R$] eligible short-side IO quantity, admitted IO quantity, and
        residual imbalance after IO admission; and
    \item[$x(o)$] the unrounded execution entitlement of order $o$.
\end{description}

\subsection{Closing Procedure}
\label{sec:close-procedure}

\begin{algorithm}[p]
\caption{Uniform-price market close}
\label{alg:market-close}
\scriptsize
\begin{algorithmic}[1]
\Require Ordinary orders $\mathcal{O}^{\mathrm{ord}}$, IO orders
    $\mathcal{O}^{\mathrm{IO}}$, and price grid $\mathcal{P}$
\Ensure A no-cross outcome, or the phase-1 and phase-2 intervals
    $\Pi_V,\Pi_\Delta$, clearing price $\pi^\star$, traded volume, and
    entitlements $x(o)$

\State Aggregate ordinary orders into $(B_i,A_i)_{i=0}^{N-1}$ by entered price
    \label{line:close-aggregate}
\For{$i=0,\ldots,N-1$}
    \State $D_i\gets\sum_{j=i}^{N-1}B_j$
        \Comment{ordinary demand eligible at $\pi_i$}
        \label{line:close-demand}
    \State $S_i\gets\sum_{j=0}^{i}A_j$
        \Comment{ordinary supply eligible at $\pi_i$}
        \label{line:close-supply}
    \State $V_i\gets\min\{D_i,S_i\}$;
        $G_i\gets S_i-D_i$;
        $\Delta_i\gets|G_i|$
        \label{line:close-point-values}
\EndFor

\State $V_{\max}\gets\max_{0\leq i<N}V_i$
    \label{line:close-vmax}
\If{$V_{\max}=0$}
    \State \Return \textsc{No-Cross}
        \label{line:close-no-cross}
\EndIf

\State $\mathcal{I}_V\gets\{i:V_i=V_{\max}\}$
    \label{line:close-plateau}
\State $c_V\gets\min\mathcal{I}_V$;
    $d_V\gets\max\mathcal{I}_V$;
    $\Pi_V\gets[\pi_{c_V},\pi_{d_V}]$
    \Comment{phase-1 interval}
    \label{line:close-volume-interval}
\State $\Delta_{\min}\gets\min\{\Delta_i:i\in\mathcal{I}_V\}$
    \label{line:close-delta-min}
\State $\mathcal{I}_{\Delta}\gets
    \{i\in\mathcal{I}_V:\Delta_i=\Delta_{\min}\}$
\State $c_\Delta\gets\min\mathcal{I}_\Delta$;
    $d_\Delta\gets\max\mathcal{I}_\Delta$;
    $\Pi_\Delta\gets[\pi_{c_\Delta},\pi_{d_\Delta}]$
    \Comment{phase-2 interval}
    \label{line:close-imbalance-interval}
\State $i^\star\gets\operatorname{TieBreak}(\mathcal{I}_\Delta)$;
    $\pi^\star\gets\pi_{i^\star}$
    \Comment{public post-processing}
    \label{line:close-price}
\State $x(o)\gets0$ for every
    $o\in\mathcal{O}^{\mathrm{ord}}\cup\mathcal{O}^{\mathrm{IO}}$

\If{$G_{i^\star}=0$}
    \ForAll{$o\in\mathcal{O}^{\mathrm{ord}}$ eligible at $\pi^\star$}
        \State $x(o)\gets q(o)$
    \EndFor
    \State \Return $(V_{\max},\Pi_V,\Delta_{\min},\Pi_\Delta,
        \pi^\star,V_{\max},x)$
        \label{line:close-balanced}
\EndIf

\If{$G_{i^\star}<0$}
    \State $\mathit{long}\gets\mathrm{BUY}$;
        $m\gets\max\{i\in\mathcal{I}_\Delta:G_i<0\}$;
        $P\gets B_m$
    \State $J\gets$ total sell-IO quantity eligible at $\pi^\star$
\Else
    \State $\mathit{long}\gets\mathrm{SELL}$;
        $m\gets\min\{i\in\mathcal{I}_\Delta:G_i>0\}$;
        $P\gets A_m$
    \State $J\gets$ total buy-IO quantity eligible at $\pi^\star$
\EndIf
\State $I\gets\min\{J,\Delta_{\min}\}$;
    $R\gets\Delta_{\min}-I$
    \Comment{admit IO on the short side}
    \label{line:close-io-admission}

\ForAll{$o\in\mathcal{O}^{\mathrm{ord}}$ on the short side and eligible at
    $\pi^\star$}
    \State $x(o)\gets q(o)$
\EndFor
\ForAll{$o\in\mathcal{O}^{\mathrm{ord}}$ on the long side and more aggressive
    than $\pi_m$}
    \State $x(o)\gets q(o)$
\EndFor
\ForAll{$o\in\mathcal{O}^{\mathrm{ord}}$ on the long side with
    $e(o)=\pi_m$}
    \State $x(o)\gets q(o)(P-R)/P$
        \label{line:close-long-ration}
\EndFor
\If{$J>0$}
    \ForAll{$o\in\mathcal{O}^{\mathrm{IO}}$ on the short side and eligible at
        $\pi^\star$}
        \State $x(o)\gets q(o)I/J$
            \label{line:close-io-ration}
    \EndFor
\EndIf
\State \Return $(V_{\max},\Pi_V,\Delta_{\min},\Pi_\Delta,
    \pi^\star,V_{\max}+I,x)$
    \label{line:close-return}
\end{algorithmic}
\end{algorithm}

All nonzero executions occur at the uniform price $\pi^\star$. The candidate
condition in Algorithm~\ref{alg:market-close} guarantees
$P\geq\Delta_{\min}>0$ in an unbalanced market. Consequently, the long-side
allocation $(P-R)/P$ lies in $[0,1]$. The IO ratio $I/J$ is used only when
$J>0$ and also lies in $[0,1]$.

The quantities $x(o)$ are unrounded entitlements. If the traded asset is
indivisible, the venue applies a published integer-rounding rule and specifies
how any residual units are distributed. For every crossing ordinary book the
candidate set $\mathcal{C}$ is nonempty, so the procedure returns either a
no-cross outcome or a clearing price and allocation.

% Notation reservations for later sections:
% - i^\star is an index; pi^\star is a price.
% - P is the at-price long-side volume and should not also name a polynomial.
% - S_i is supply; use a different symbol for the lookup running sum.
% - N is the number of real price ticks; reserve n for a padded domain size.
